\section{Introduction}\label{sec:intro}
% =============================================================================================
% Owner: front writer. Plan and claims: paper/OUTLINE.md, Section 3. Notation: notation.md. No
% tables or figures in this section. Every number is taken from a table that a result section
% inputs: tab:e2-size, tab:e2-power, tab:e2dep-power (fig:e2dep-power), tab:e2-llr
% (manual_e2_llr_main), tab:e2-relative-boundary (manual_e2_relative_boundary_main),
% tab:e2-dependent-size, tab:e1-size-sim (manual_e1_size_sim_main), tab:e1-nonstat
% (e1_nonstat), tab:e3-stylized, tab:e3-power (manual_e3_power_main), tab:e3-pooled-tests,
% tab:e3-verdicts, tab:e3-agreement, tab:e4-correlation, tab:e4-utility, tab:e4-utility-agg,
% tab:e4-tests, tab:e4-link. The 45% figure is the theory.md pilot (Section 11.5), quoted as a
% pilot only. The count "all 14 configurations" (critic ranking of the likelihood models in the
% E5 grid) was checked on results/e5/e5_sensitivity.csv (median absolute statistic).
% Post-hoc architecture ablation: tab:e5-arch (e5_arch_sim.tex; 0.98 / 0.10 / 1.00 for t10 vs t4,
% n=500); "pilot on four E2 designs" and "D2 and seed combination added on the day of the runs"
% from the docs/DESIGN.md change log; "six tickers" from results/e3_quick/meta (development runs).
% "Barely reproducible" (E4 critic statistic): seed-0 vs seed-1 Spearman over the 30 universes
% within each generator, -0.03 to 0.30 (tab:manual-e4-reliability, from results/e4/e4_tests.csv).
% Pooled vs replica-wise variance ratios (MV Gaussian GMV 1.038 / 1.004): tab:e4-utility-agg.
% Fixed realised-volatility statistic (0.69-0.94 of the critic's z_hac, same ranking of GBM,
% Merton, GARCH-t):
% tab:e3-rv-baseline. Indep-cell oversize vs standard error: tab:e1-sim-se (sec:sim-dependent).
% =============================================================================================

Synthetic return data, simulated from a fitted model or from a trained generative network, now
serve several purposes in finance. They are used to back-test trading strategies on many paths
instead of the single historical one, which a strategy selected in-sample tends to overfit
\citep{bailey2017probability}; to share data that cannot be disclosed for confidentiality or
privacy reasons \citep{assefa2020generating}; to train hedging strategies
\citep{buehler2019deep,kondratyev2019market}; and to augment short samples, for instance when
portfolio inputs are estimated \citep{buehler2020data,cetingoz2026synthetic}. The generators
range from classical parametric models (geometric Brownian motion, jump-diffusions, GARCH,
stochastic volatility) to generative networks
\citep{wiese2020quant,takahashi2019modeling,buehler2020data,ni2021sig}. Each of these uses
assumes an answer to the question in our title: is the synthetic data good enough?

The question is usually answered informally. Most papers compare stylized facts of real and
generated paths, such as heavy tails, volatility clustering and leverage
\citep{cont2001empirical}, visually or through summary statistics, and evaluation practice varies
widely \citep{stenger2024evaluation}. Discriminator-based scores, such as the held-out error of a
classifier trained to tell real from synthetic series, are used without a null distribution
\citep{yoon2019time,francopedroso2018ets,storchan2021learning}. Without a null distribution and a
known power, a low score cannot be told apart from a weak classifier. Formal two-sample tests have
been used
to evaluate market generators \citep{buehler2020data,chevyrev2022signature}, but without an
analysis of their size and power on financial data. Different evaluation criteria for generative
models can also disagree sharply \citep{theis2016note}, so a criterion has to answer a precise
statistical question and, ideally, say something about the intended use.

We treat the question as a hypothesis test. A yes/no answer that a statistician can defend needs
five things.
\begin{enumerate}
  \item[(i)] \emph{Level control on dependent data.} Real return windows are serially dependent
  through volatility clustering and cross-sectionally dependent through common market shocks, so
  permutation p-values, which assume exchangeability, are not valid for them. With windows of many
  stocks pooled and the training and evaluation halves split by ticker, a pilot simulation of a
  stylised one-factor panel rejected a true null in about 45\% of replications at nominal 5\%. In
  the full placebo of E1, a factor panel fitted to our data and run through the complete pipeline,
  the same protocol rejected in all 200 replications (\cref{tab:e1-size-sim}).
  \item[(ii)] \emph{Power against the deviations that matter in finance}, quantified before the
  test is applied to real data. Without known power, a non-rejection says nothing about quality.
  \item[(iii)] \emph{A null that respects non-stationarity.} A generator $G$ is calibrated on the
  real past $\Rpast$ and judged against later real data $\Rnew$, and it inherits features of the
  past that need not persist. In our data every generator fitted to 2018--2022 is more volatile
  than the 2023--2025 returns (panel median daily standard deviation between 1.89\% and 2.13\%,
  against 1.62\%; \cref{tab:e3-stylized}). The absolute null $\Law(G)=\Law(\Rnew)$ asks the
  generator to anticipate such changes. Users who back-test or manage risk need exactly that, and
  we report this test, but its rejection does not separate a poor generator from a change of
  regime after the calibration period. As in forecast comparisons against a benchmark model
  \citep{diebold1995comparing,giacomini2006tests}, we therefore also compare $G$ with the naive
  benchmark of replaying history, that is, historical simulation from $\Rpast$: is $G$ further from
  $\Rnew$ than $\Rpast$ is? A rejection then means worse than that benchmark, not merely different
  from the future.
  \item[(iv)] \emph{A link to likelihood-based evaluation.} Econometricians trust it, and where a
  likelihood exists it provides a benchmark verdict and a power envelope.
  \item[(v)] \emph{Evidence on use}: whether the verdict predicts how useful the data are for a
  given decision.
\end{enumerate}

This paper builds such a test around the critic of a Wasserstein GAN, trained alone, with no
generator. A critic is a neural-network test function chosen to maximise a mean difference
between two samples under a Lipschitz penalty \citep{arjovsky2017wasserstein,gulrajani2017improved}.
Let $X_1,\dots,X_n$ be windows of $L$ consecutive daily returns of real data, with law $P$, and let
$Y_1,\dots,Y_m$ be windows simulated from $G$, with law $Q$. We split both samples into a training
half and an evaluation half and train the critic $\fhat$ on the training halves. The statistic for
the absolute null $P=Q$ (\cref{def:abs-null}) is the held-out mean difference
\begin{equation*}
  T=\Pev\fhat-\Qev\fhat=\frac{1}{n_e}\sum_{i\in\Iev}\fhat(X_i)-\frac{1}{m_e}\sum_{j\in\Jev}\fhat(Y_j).
\end{equation*}
Given the training half, $T$ is unbiased for $P\fhat-Q\fhat$. This quantity is bounded by the
Lipschitz constant of the network times the Wasserstein-1 distance between the standardised
window laws (\cref{prop:unbiased}). We read $T$ as a learned moment that can be evaluated window
by window (\cref{sec:emp-witness}), not as an estimate of $\W$.

For an econometrician the construction is a learned, sample-split relative of the consistent
specification tests of \citet{bierens1990consistent} and \citet{stinchcombe1998consistent}: they
use a whole class of test functions, whereas we select one on a training split and test with it
on the other. It is also close in spirit to the simulation-based specification tests of
\citet{corradi2005bootstrap}, with one difference that matters. Their null is that the model class
is correctly specified. Ours is that the fitted generator, with parameters estimated on a
calibration period and then held fixed, reproduces the window law of a later evaluation period
(\cref{sec:lit-econometrics}). None of the building blocks is new on its own, and
\cref{sec:lit-summary} lists what we take from earlier work.

\paragraph{Contributions.}
\begin{enumerate}
  \item \emph{A critic test with known level.} For independent windows, the permutation test with
  the critic held fixed has exact finite-sample level whatever the training did
  (\cref{prop:exact-level}). For dependent panels we split by contiguous time blocks rather than
  by ticker, so that training and evaluation windows cover disjoint dates, and we studentise $T$
  with a Driscoll--Kraay variance \citep{driscoll1998consistent}. The resulting p-value is
  asymptotically valid under mixing conditions for a fixed critic whose evaluation data are
  independent of its training data (\cref{prop:hac}). A ticker split on common dates would
  instead couple the two halves through the realised market path, and the statistic is then
  biased under the null (\cref{prop:ticker-bias}). Our time-block split interleaves training and
  evaluation blocks without a gap, so it does not deliver that independence
  (\cref{lem:decoupling}), and we measure the size of the complete real-data pipeline in a
  simulated placebo panel calibrated to the data: at nominal 5\% it rejects a true null at 0.085,
  against 1.000 for the ticker split with permutation p-values, but it over-rejects when the
  synthetic tickers are independent, as they are for our univariate generators
  (\cref{sec:emp-placebo,sec:sim-dependent}).

  \item \emph{An anchored relative test and a three-valued verdict.} Relative similarity tests
  compare two candidate models with one reference sample \citep{bounliphone2016test}. In ours one
  candidate is the real past, so the null states that $G$ is at least as close to the future as
  the past is. A single critic is trained to separate $\Rnew$ from an equally weighted pool of
  synthetic and past training windows. The statistic $\Danc$ is its mean on held-out past windows
  minus its mean on held-out synthetic windows (\cref{eq:anchored-stat}). At the boundary
  $\Law(G)=\Law(\Rpast)$ the two held-out samples are exchangeable given the critic, so for
  independent windows the permutation test is exact for every critic (\cref{prop:anchored}). The
  obvious alternative, the difference of two separately trained critic statistics, has no such
  guarantee: at that boundary its size is at least the nominal level, and for fixed critics and
  growing evaluation sets it tends to $1/2$ (\cref{prop:boundary-size}). A generator is \BAD{} if
  the anchored test rejects. It is \GOOD{} if the test does not reject and a power check shows
  that in simulated reference designs the same pipeline detects a synthetic volatility 20\% too
  high with probability at least $0.8$. It is \INCONCL{} otherwise
  (\cref{def:power-check,def:verdict,prop:verdict}). On real data, where windows are dependent and
  the past is itself non-stationary, the verdict combines Driscoll--Kraay p-values over several
  synthetic seeds; the exactness results above do not cover this rule, so the power check also
  estimates its size at the boundary, in a stationary design and in one that keeps the real
  volatility path of the past (\cref{sec:method-decision}). The check was fixed before the full
  real-data runs; the second design and the seed combination were added on the day of those runs,
  after development runs on six tickers (\cref{sec:emp-power}).

  \item \emph{Calibration against ground truth and likelihood-ratio oracles.} On independent
  simulated windows we measure the size and power of the \WC{} test and of five competitors: MMD
  \citep{gretton2012kernel}, a signature MMD \citep{chevyrev2022signature}, the classifier
  two-sample test \citep[C2ST;][]{lopezpaz2017revisiting}, a learned log-likelihood-ratio
  statistic \citep[\LRC{};][]{cheng2022classification} and the Kolmogorov--Smirnov (KS) test on
  pooled returns. The benchmark is the Neyman--Pearson oracle, which no two-sample test can beat
  (\cref{prop:oracle-bound}). The theory predicts a large gap. With a 50/50 split and equal sample
  sizes, the local noncentrality of any test based on the held-out mean difference of a fixed
critic is at most half the oracle's
  (\cref{prop:local-power}; proved for a fixed critic, sketched for a trained one). The comparison
  states where simpler tests win, and it includes alternatives that change only the dependence
  within a window (\cref{sec:sim}).

  \item \emph{Likelihood benchmarks on S\&P~500 data.} Five univariate generators (GBM, Merton,
  \GARCHt{}, Heston and the stationary bootstrap) are fitted to 100 stocks on 2018--2022 and
  evaluated on 2023-01-03 to 2025-06-16. Three of them have a tractable likelihood (GBM, Merton,
  \GARCHt{}). For these we compare the verdicts with out-of-sample predictive log-likelihood,
  probability-integral-transform tests and Amisano--Giacomini comparisons
  \citep{diebold1998evaluating,berkowitz2001testing,amisano2007comparing}. Heston has no
  closed-form likelihood \citep{aitsahalia2007maximum} and the bootstrap no parametric one; these
  are the cases that call for a likelihood-free test (\cref{sec:emp}). The critic ranks GBM,
  Merton and \GARCHt{} exactly as the predictive log-score does; GBM, Merton and Heston are \BAD{}
  and \GARCHt{} is \INCONCL{} (\cref{sec:emp-verdicts}).

  \item \emph{A link to Markowitz utility.} We show that critic distance and portfolio utility can
  disagree in both directions (\cref{sec:mk-theory}). The utility of the global minimum-variance
  (GMV) portfolio depends on the covariance matrix only up to scale, so a generator whose only
  error is a scalar factor loses nothing for GMV, yet a consistent critic test rejects it
  (\cref{ex:scale-blind}). Conversely, an error in a low-variance direction, which GMV weights
  heavily, can be nearly invisible in $\W$ (\cref{ex:spread}). Empirically, we build GMV and
  mean-variance portfolios from the output of five multivariate generators, evaluate them on the
  real evaluation period against the same portfolios estimated on the real calibration period,
  and relate the utility losses to multivariate critic and MMD statistics. The statistics do not
  predict the losses; the distance between the synthetic and the past average correlation does,
  with a negative sign, consistent with the fall of the real correlations, a sign that could
  reverse in a period of rising correlations (\cref{sec:markowitz}).
\end{enumerate}

\paragraph{Main findings.}
We call \emph{window tests} the tests that compare the laws of whole windows (\WC{}, \LRC{}, C2ST,
MMD, \SigMMD{}), as opposed to KS on the $nL$ pooled returns. With independent windows every
window test holds its level: in each family and at both sample sizes, the Wilson 95\% interval of
its rejection rate at nominal 5\% contains 0.05 (\cref{tab:e2-size}). KS, which treats the pooled
returns as independent, is not valid when returns within a window are dependent: returns within a
GARCH window share their conditional variance, and KS rejects a true null at 0.110 [0.074, 0.161]
with $n=500$ windows and at 0.080 [0.050, 0.126] with $n=2000$ (Wilson 95\% intervals in
brackets); in the other three families its intervals contain 0.05.
Against deviations of the one-dimensional marginal (volatility level, tails, jumps), KS is
nevertheless far more powerful than every window test. At $n=2000$ it rejects Gaussian against
Student-$t_4$ returns at matched variance, and a variance-matched GBM against a Merton
jump-diffusion, in every replication; the best window test, the \WC{} test, rejects at 0.67 and
0.20. Against volatility clustering at matched variance the window tests have power: with $n=2000$
and GARCH persistence 0.95 against 0.98, they reject at 0.79 (\WC{}), 0.93 (MMD), 0.86
(\SigMMD{}), 0.55 (C2ST) and 0.30 (\LRC{}) (\cref{tab:e2-power,fig:e2-power}). KS rejects even more
often there, but it is oversized and its power comes from the thinner tails that lower persistence
gives the marginal. When the marginal is held exactly fixed, KS has no power and the window tests
detect the deviation: the \WC{} test rejects at 0.46, 0.94 and 1.00 as the share of windows made
of independent returns rises from one half to one, below MMD and \SigMMD{} at one half (0.70)
and comparable to them at the larger shares, and above C2ST and the \LRC{}
(\cref{fig:e2dep-power}). No window test dominates. A kernel test is the most
powerful window test, or tied for it, at every alternative of volatility level and of
clustering, while the \WC{} test is the most powerful against the lightest-tailed alternatives to
the Student-$t_4$ law ($t_{30}$ and Gaussian). At the smallest deviation of each family with
$n=500$, every window test stays far below the oracle, which rejects in at least 0.98 of the
replications at every alternative of the four families; against a quarter
or a half of GBM windows in a Merton sample no window test rejects more often than 0.08. The
learned scores are tests, not likelihood estimators: their Spearman correlation with the exact
log-likelihood ratio lies between 0.024 and 0.556 (\cref{tab:e2-llr}). At the relative boundary
$\Law(G)=\Law(\Rpast)$, with $n=500$, the anchored test rejects at 0.065 [0.038, 0.108] and 0.055
[0.031, 0.096] in two designs, and the two-critic test at 0.040 and 0.065, so the excess that
\cref{prop:boundary-size} allows is not visible at this sample size. A generator that resamples
the past makes the anchored test reject at 0.965 with a random split, against 0.010 when the split
respects the source window (\cref{tab:e2-relative-boundary}). With consecutive windows of one GARCH
path, the random-split permutation test is oversized (0.145 for the \WC{} test, 0.290 for MMD),
and the Driscoll--Kraay p-value after a time-block split and the wild bootstrap bring the rates to
0.070 and 0.075 (\cref{tab:e2-dependent-size}).

In a simulated factor panel with the dimensions of our data, the time-block Driscoll--Kraay
protocol rejects a true null at 0.085 [0.054, 0.132] for the \WC{} test (0.050 against a fixed-$b$
reference) and the wild-bootstrap MMD test at 0.045. The ticker split rejects at 0.990 even with
Driscoll--Kraay standard errors, the permutation p-value after a time-block split at 0.410 and KS
at 0.970 (\cref{tab:e1-size-sim}). When the synthetic tickers are drawn independently, as the
univariate generators draw them, the learned tests over-reject (0.270 for the absolute \WC{} test,
0.155 for the anchored test). The statistic itself is no more biased than when the synthetic
panel has a factor, but it is negatively correlated with its Driscoll--Kraay standard error, and
the rejections concentrate in the replications with a small standard error
(\cref{tab:e1-sim-se}); \cref{cor:block-bias} gives a model-based explanation, which we have not
tested. Tested against each other with the same protocol, the real \TRAIN{} and \TEST{} windows are not rejected as
different ($\pmed=0.452$ for the \WC{} test; \cref{tab:e1-nonstat}), but with few evaluation dates
(\cref{rem:small-T}) this non-rejection says little about whether the two periods are similar. The
relative null is motivated by the volatility level that the generators inherit, not by a
rejection of the past. The power check shows the cost of non-stationarity. In a stationary
reference design the verdict rule detects a 20\% volatility error in all 100 replications and
never rejects at the boundary; with the real volatility path of 2018--2022 its power is 0.45 [0.36, 0.55] and its boundary size 0.08 [0.04,
0.15] (\cref{tab:e3-power}), so on this panel and split no generator can be \GOOD. The absolute
\WC{} test rejects all five generators (\cref{tab:e3-pooled-tests}). The anchored test rejects GBM
($\pmed<0.001$), Merton (0.005) and Heston (0.017), which, as calibrated on 2018--2022, are \BAD{},
Heston borderline. It does not reject \GARCHt{} (0.432), which is distinguishable from the future
but not detectably worse than history, and \INCONCL{}; the stationary bootstrap, which copies the
past, receives no verdict (\cref{tab:e3-verdicts}). The witness shows that the critic mainly
detects a volatility level: every generator is too volatile for 2023--2025
(\cref{sec:emp-witness}). Each verdict therefore concerns a model together with its calibration
on a period that contains the 2020 and 2022 turbulence, not the model class alone. The critic ranks
\GARCHt{}, Merton and GBM exactly as the predictive log-score does (Kendall $\tau=1$), whereas its
per-ticker rejections agree little with Berkowitz tests (Cohen $\kappa=0.11$), which check one-day
calibration rather than simulated windows (\cref{tab:e3-agreement}). In the sensitivity analysis
this ranking holds in all 14 configurations, but the \BAD{} verdicts of Merton and Heston do not,
and with an alternative split every parametric generator is \INCONCL{} (\cref{sec:robust}).
Fitting is itself informative. Simulated with the quadratic-exponential scheme of \citet{andersen_2008_simple},
Heston needs a bounded volatility of variance to avoid days with zero variance, and it then falls
well short of the kurtosis of the data (\cref{tab:e3-heston-moments}).

In the Markowitz study, the average pairwise correlation of real returns fell from 0.419 in
\TRAIN{} to 0.268 in \TEST{} (\cref{tab:e4-correlation}). Four of the five generators give
minimum-variance portfolios whose \TEST{} variance is within 6\% of that of the portfolio trained
on real data (replica-wise variance ratios). With pooled moments, CCC-GARCH-$t$ and MV Student-$t$
beat it (variance ratios 0.869 and 0.896); their average correlations are the furthest below the
past one. Factor GARCH-$t$ loses most (1.466 with pooled moments, 1.436 replica-wise) although its
average correlation is close to the past one (\cref{tab:e4-utility,tab:e4-utility-agg}). In this
evaluation period test results do not predict these losses. MV Gaussian is rejected in all 30
universes after size adjustment (first seed), yet its minimum-variance portfolios have \TEST{}
variance ratios close to one (1.038 replica-wise, 1.004 with pooled moments). Within a generator
the critic statistic is barely reproducible across synthetic seeds (Spearman correlation between
the two seeds from $-0.03$ to 0.30; \cref{tab:manual-e4-reliability}), so its near-zero
correlation with the minimum-variance loss (0.01) is weak evidence. The losses are point
estimates over one \TEST{} period, without inference over dates. The one consistent correlate is the distance between the synthetic and the past average correlation, with
a negative sign ($-0.36$ within generators; \cref{tab:e4-tests,tab:e4-link}), which reflects the
fall in correlations and need not hold in a period in which they rise.

\paragraph{What we do not claim.}
The \WC{} test is one well-calibrated and interpretable tool within a validated protocol, not a
uniformly most powerful test. KS is more powerful against marginal deviations, and a kernel test
is at least as powerful as the \WC{} test at every alternative of volatility level and of
clustering in our simulations. The power of the critic depends on its optimiser settings and its
architecture (\cref{sec:robust}). In a post-hoc ablation, whose architecture was chosen after the
real-data results of the pre-specified perceptron were known and after a pilot on four of the
simulated designs, a time-pooled convolutional critic keeps the level and comes close to KS
against tails and jumps: with $n=500$ it rejects $t_{10}$ against $t_4$ windows at 0.98, where the
perceptron rejects at 0.10 and KS at 1.00 (\cref{tab:e5-arch}). The comparison with KS is thus a
statement about the pre-specified network, whose verdicts are the ones we report. We do not
compare with power-optimised deep-kernel tests \citep{liu2020learning}, and on real data, where
the critic behaves almost like a realised-volatility statistic (\cref{sec:emp-witness}), a fixed
log realised-volatility statistic on the same seed-0 windows ranks the three likelihood models
in the same order and reaches between 0.69 and 0.94 of the critic's $\zhac$
(\cref{tab:e3-rv-baseline}); we did not run it inside the verdict rule. The protocol contributes exact or asymptotically valid inference on
dependent panels, the relative null, the link with likelihood and an extension to multivariate
windows, where pooled KS sees only the marginal; the size of its dependent-data version is
measured, not proved, and it is not nominal in every design. A \GOOD{} verdict means that a
procedure with demonstrated power against a deviation of pre-specified size did not detect one. It
never means that the generator equals the truth, and its error control rests on a premise that
cannot be verified (\cref{prop:verdict}). The relative null does not penalise a generator error
that happens to point towards the change of regime, and the power check measures power only
against a volatility 20\% too high, an error that on our data points away from the calmer
evaluation period. A rejection is conditional on the fitted parameters, so
it does not separate misspecification from estimation error or from parameter drift between the
two periods; the relative test is designed to absorb the drift. Under non-stationarity a verdict
concerns one pair of periods. Univariate verdicts say nothing about the cross-sectional
dependence that drives portfolio risk (\cref{ex:corr-blind}), and whether synthetic data are good
for a task has to be checked on that task. Finally, the panel consists of the S\&P~500
constituents at the end of the sample. This is a selection on outcomes of the evaluation period
that goes beyond survivorship bias \citep{brown1992survivorship,shumway1997delisting}, and it may
make the real evaluation panel calmer than the one a forecaster faced at the end of 2022.

\paragraph{Outline.}
\Cref{sec:lit} reviews related work and \cref{sec:method} defines the procedure. \Cref{sec:theory}
states the theory: exact level, consistency, the likelihood link, the relative test, inference
under dependence and what the verdict controls. \Cref{sec:sim} reports the simulation evidence,
\cref{sec:data} describes the data and the generators, \cref{sec:emp} gives the S\&P~500 results
and \cref{sec:markowitz} studies Markowitz utility. \Cref{sec:robust} checks robustness and
\cref{sec:concl} concludes. Proofs are in \cref{app:proofs}, implementation details in
\cref{app:impl} and further results in \cref{app:sim,app:data,app:emp,app:mk,app:robust}.
